
Deep neural networks trained with empirical risk minimization (ERM) on datasets with spurious correlations systematically underperform on minority groups. On Waterbirds, where 95\% of waterbirds appear on water backgrounds, ERM achieves approximately 97\% average accuracy but only approximately 75\% on the minority group (waterbirds on land). This accuracy gap has motivated interventions including Group DRO, last-layer retraining, and various reweighting schemes.

These methods address the symptom of poor minority accuracy without explaining the mechanism. Recent work established that minority groups occupy sharper regions of the loss landscape, with sharpness ratios SR > 1.0. Whether this curvature asymmetry is intrinsic to the data or emerges from training dynamics remained an open question.

This work investigates the Differential Convergence Rate (DCR) hypothesis: majority groups, having more samples, converge faster during training. Their gradient norms decay earlier, smoothing curvature in majority-relevant parameter directions. Minority-relevant directions, receiving fewer updates, remain sharp.

\textbf{Contributions:}

\begin{enumerate}
\item Establishment that SR $\approx$ 1.0 at random initialization ($SR_0$ = 0.9999, SEM = 0.0017), indicating curvature asymmetry is training-induced rather than data-intrinsic.
\end{enumerate}

\begin{enumerate}
\item Demonstration that gradient ratio decay temporally precedes sharpness ratio divergence by $\tau$ = 3.67 epochs (95\% CI [3.0, 4.0]), consistent with a causal relationship.
\end{enumerate}

\begin{enumerate}
\item Identification of a 3-4 epoch diagnostic window between gradient convergence and curvature divergence.
\end{enumerate}

